\documentclass[11pt,a4paper]{article}

\usepackage[utf8]{inputenc} % allow utf-8 input
\usepackage[T1]{fontenc}    % use 8-bit T1 fonts
\usepackage{hyperref}       % hyperlinks
\usepackage{url}            % simple URL typesetting
\usepackage{booktabs}       % professional-quality tables
\usepackage{amsfonts}       % blackboard math symbols
\usepackage{nicefrac}       % compact symbols for 1/2, etc.
\usepackage{microtype}      % microtypography
\usepackage{xcolor}         % colors
\usepackage{graphicx}       % graphics
\usepackage[normalem]{ulem} % strikethrough
\title{2nd Sample Document Title}

\begin{document}

\maketitle

This is an abstract.

\section{Section 1: Testing nestedlists}

\begin{itemize}
\item First item
  \begin{itemize}
\item Nested item 1
\item Nested item 2
  \end{itemize}
\item Second item
  \begin{enumerate}
\item Nested ordered item 1
\item Nested ordered item 2
    \begin{itemize}
\item Deeper nested unordered item
    \end{itemize}
  \end{enumerate}
\item Third item
  \begin{enumerate}
\item Nested ordered item 1
\item Nested ordered item 2
    \begin{itemize}
\item Deeper nested unordered item
    \end{itemize}
\item Nested ordered item 2
  \end{enumerate}
\end{itemize}

\section{Section 2}

bla bla

bla bla bla bli bla ble

\section{Section 3: test image}

\section{Section 4: test tables}

\begin{table}[h]
\begin{tabular}{|l|l|l|}
\hline
Header 1 & Header 2 &  \\ \hline
Value 1 & Value 2 &  \\ \hline
Value 3 & Value 4 &  \\ \hline
\end{tabular}
\end{table}

Caption for the table 1

\begin{table}[h]
\begin{tabular}{|l|l|}
\hline
Header 1 & Header 2 \\ \hline
Value 1 & Value 2 \\ \hline
Value 3 & Value 4 \\ \hline
\end{tabular}
\end{table}

Caption for the table 2

\begin{table}[h]
\begin{tabular}{|l|l|l|}
\hline
Column 1 Heading & Column 2 Heading & Column 3 Heading \\ \hline
Cell 1 & Cell 2 & Cell 3 \\ \hline
Cell 4 & Cell 5 colspan=2 & Cell spans two columns \\ \hline
\end{tabular}
\end{table}

Caption for the table 3

\begin{table}[h]
\begin{tabular}{|l|l|l|}
\hline
Column 1 Heading & Column 2 Heading & Column 3 Heading \\ \hline
Rowspan=2 & Cell 2 & Cell 3 \\ \hline
 & Cell 5 & Cell 6 \\ \hline
\end{tabular}
\end{table}

Caption for the table 4

\begin{table}[h]
\begin{tabular}{|l|l|l|l|}
\hline
Col 1 & Col 2 & Col 3 & Col 4 \\ \hline
Rowspan=2.Colspan=2 & Cell spanning 2 rows and 2 columns & Col 3 & Col 4 \\ \hline
 &  & Col 3 & Col 4 \\ \hline
Col 1 & Col 2 & Col 3 & Col 4 \\ \hline
\end{tabular}
\end{table}

Table 5 with multiple empty cells

\begin{table}[h]
\begin{tabular}{|l|l|l|l|}
\hline
Column 1 Heading & Column 2 Heading & Column 3 Heading &  \\ \hline
Cell 1 &  & Cell 3 &  \\ \hline
Cell 4 &  &  &  \\ \hline
Cell 7 &  &  &  \\ \hline
Cell 10 & Cell 11 & Cell 12 &  \\ \hline
 & Cell 14 & Cell 15 &  \\ \hline
 &  &  &  \\ \hline
Cell 19 & Cell 20 & Cell 21 &  \\ \hline
\end{tabular}
\end{table}

\subsubsection{SubSubSection 2.1.1}

\begin{figure}[h]
% image
\end{figure}

\begin{figure}[h]
% image
\caption{An example caption for the image}
\end{figure}

\end{document}